\problema{All Nodes Distance K in Binary Tree}{863}{src/02-arboles/863-all-nodes-distance-k-in-binary-tree.cpp}{Media}
\complejidad{O(N)}{O(N)}

Se nos da un árbol binario, un nodo objetivo y un número $K$, y hay que devolver todos los nodos que están exactamente a $K$ pasos de distancia del objetivo.

\paragraph{La idea, con un ejemplo de la vida real.} En un \emph{árbol binario} cada \emph{nodo} (cada ``persona'') solo sabe quiénes cuelgan de él (sus hijos izquierdo y derecho), pero no sabe quién está arriba (su \emph{padre}). Es como un juego de las escondidas donde solo puedes ver hacia abajo. El problema es que los nodos que buscan pueden estar hacia abajo del objetivo, pero también hacia arriba y luego bajando por otra rama, y para llegar a esos hay que poder subir.

La solución es como tratar el árbol igual que un mapa de amistades: si en vez de ``padre e hijos'' pensamos que cada nodo simplemente tiene \emph{vecinos}, entonces cada nodo tiene hasta tres vecinos: su hijo izquierdo, su hijo derecho y su padre. Ese mapa de vecinos (donde las conexiones se pueden recorrer en los dos sentidos) se llama \emph{grafo}, y ahí sí podemos caminar libremente en todas direcciones.

\paragraph{Cómo lo resuelve el código, paso a paso.} Lo hacemos en dos etapas.
\begin{enumerate}
  \item \textbf{Averiguar quién es el padre de cada quien.} Recorremos el árbol una vez bajando por las ramas y guardamos en un diccionario (un \emph{mapa hash}, que es como una agenda que empareja cada nodo con un dato) quién es el padre de cada nodo. Así cada nodo ya conoce a sus tres vecinos.
  \item \textbf{Expandirnos como una gota de tinta en agua.} Partimos del objetivo y avanzamos por niveles, tocando primero a los vecinos que están a 1 paso, luego a los que están a 2 pasos, y así. Esta forma de explorar ``en círculos que crecen'' se llama \emph{BFS} (búsqueda en anchura). Usamos una \emph{cola} (una fila, como la del supermercado: el primero que entra es el primero que sale) y un \emph{conjunto} de visitados para no regresar por donde ya pasamos. Cuando hemos avanzado exactamente $K$ niveles, todos los nodos que quedaron en la fila son justo los que están a distancia $K$.
\end{enumerate}

\lstinputlisting{src/02-arboles/863-all-nodes-distance-k-in-binary-tree.cpp}

\paragraph{¿Qué significa esa complejidad?} $N$ es la cantidad de nodos. Tanto el tiempo como la memoria son $O(N)$: en el peor caso visitamos cada nodo una sola vez, y el mapa de padres más la cola guardan a lo mucho un dato por nodo. No se puede hacer más rápido, porque para estar seguros hay que poder mirar todo el árbol.

\paragraph{Consejos para la entrevista (Amazon).} La idea clave que el entrevistador quiere oír es ``convertir el árbol en un grafo para poder subir a los padres''. Marcar los nodos como visitados es imprescindible: si se te olvida, la búsqueda rebota entre un nodo y su padre para siempre. Menciona también los casos raros: si $K$ es 0 la respuesta es solo el propio objetivo, y si $K$ es más grande que el árbol, la respuesta queda vacía.
